% Chapter 4: Numerical Computation

\chapter{Numerical Computation}
\label{chap:numerical-computation}

This chapter covers numerical methods and computational considerations essential for implementing deep learning algorithms. Topics include gradient-based optimization, numerical stability, and conditioning.

\section*{Learning Objectives}

After studying this chapter, you will be able to:

\begin{itemize}
    \item \textbf{Understand numerical precision issues}: Recognize how finite precision arithmetic affects deep learning computations and implement solutions for overflow, underflow, and numerical instability.
    
    \item \textbf{Master gradient-based optimization}: Apply gradient descent and understand the role of Jacobian and Hessian matrices in optimization landscapes, including critical points and saddle points.
    
    \item \textbf{Handle constrained optimization}: Use Lagrange multipliers and KKT conditions to solve optimization problems with constraints, relevant for regularization and fairness in deep learning.
    
    \item \textbf{Assess numerical stability}: Calculate condition numbers, identify ill-conditioned problems, and implement gradient checking and other numerical verification techniques.
    
    \item \textbf{Apply practical numerical techniques}: Use log-sum-exp tricks, mixed precision training, and other numerical stability methods in real deep learning implementations.
    
    \item \textbf{Debug numerical issues}: Recognize common numerical problems in deep learning and apply appropriate solutions for robust model training.
\end{itemize}

\input{chapters/chap04-sec01}
\input{chapters/chap04-sec02}
\input{chapters/chap04-sec03}
\input{chapters/chap04-sec04}

\section{Problems}
\label{sec:problems}

This section contains exercises to reinforce your understanding of numerical computation concepts. Problems are categorized by difficulty level.

\subsection{Easy Problems (6 problems)}

\begin{problem}[Floating-Point Basics]
\label{prob:float-basics}
Consider a hypothetical 4-bit floating-point system with 1 sign bit, 2 exponent bits, and 1 mantissa bit. What is the smallest positive number that can be represented? What is the largest number?

\textbf{Hint:} Use the IEEE 754 format: $(-1)^s \times 2^{e-b} \times (1 + m)$ where $s$ is the sign bit, $e$ is the exponent, $b$ is the bias, and $m$ is the mantissa.
\end{problem}

\begin{problem}[Softmax Stability]
\label{prob:softmax-stability}
Compute the softmax of $\vect{x} = [1000, 1001, 1002]$ using both the naive approach and the numerically stable approach. Show your work step by step.

\textbf{Hint:} Use the identity $\text{softmax}(\vect{x}) = \text{softmax}(\vect{x} - c)$ where $c = \max_i x_i$.
\end{problem}

\begin{problem}[Gradient Computation]
\label{prob:gradient-computation}
Compute the gradient of $f(x, y) = x^2 + 3xy + y^2$ at the point $(2, 1)$.

\textbf{Hint:} The gradient is $\nabla f = \left[\frac{\partial f}{\partial x}, \frac{\partial f}{\partial y}\right]$.
\end{problem}

\begin{problem}[Critical Points]
\label{prob:critical-points}
Find all critical points of $f(x, y) = x^3 - 3x + y^2 - 2y$ and classify them as local minima, local maxima, or saddle points.

\textbf{Hint:} Set $\nabla f = \boldsymbol{0}$ and use the Hessian matrix to classify.
\end{problem}

\begin{problem}[Condition Number]
\label{prob:condition-number}
Calculate the condition number of the matrix $\mat{A} = \begin{bmatrix} 1 & 0.5 \\ 0.5 & 1 \end{bmatrix}$.

\textbf{Hint:} For a 2×2 matrix, $\kappa(\mat{A}) = \frac{\lambda_{\max}}{\lambda_{\min}}$ where $\lambda_{\max}$ and $\lambda_{\min}$ are the eigenvalues.
\end{problem}

\begin{problem}[Gradient Checking]
\label{prob:gradient-checking}
Implement gradient checking for the function $f(x) = x^3 + 2x^2 + x$ at $x = 1$ using finite differences with $\epsilon = 0.001$.

\textbf{Hint:} Use the formula $\frac{\partial f}{\partial x} \approx \frac{f(x+\epsilon) - f(x-\epsilon)}{2\epsilon}$.
\end{problem}

\subsection{Medium Problems (5 problems)}

\begin{problem}[Log-Sum-Exp Trick]
\label{prob:log-sum-exp}
Derive the log-sum-exp trick: $\log\left(\sum_{i=1}^n \exp(x_i)\right) = c + \log\left(\sum_{i=1}^n \exp(x_i - c)\right)$ where $c = \max_i x_i$.

\textbf{Hint:} Start by factoring out $\exp(c)$ from the sum.
\end{problem}

\begin{problem}[Hessian Eigenvalues]
\label{prob:hessian-eigenvalues}
For the function $f(x, y) = x^4 + y^4 - 4xy$, find the Hessian matrix and its eigenvalues at the critical point $(0, 0)$. What does this tell you about the nature of this critical point?

\textbf{Hint:} Compute $\mat{H} = \begin{bmatrix} \frac{\partial^2 f}{\partial x^2} & \frac{\partial^2 f}{\partial x \partial y} \\ \frac{\partial^2 f}{\partial y \partial x} & \frac{\partial^2 f}{\partial y^2} \end{bmatrix}$.
\end{problem}

\begin{problem}[Lagrange Multipliers]
\label{prob:lagrange-multipliers}
Find the maximum value of $f(x, y) = xy$ subject to the constraint $x^2 + y^2 = 1$ using Lagrange multipliers.

\textbf{Hint:} Set up the Lagrangian $\mathcal{L}(x, y, \lambda) = xy + \lambda(1 - x^2 - y^2)$ and solve the system of equations.
\end{problem}

\begin{problem}[Numerical Differentiation]
\label{prob:numerical-differentiation}
Compare the accuracy of forward difference, backward difference, and central difference formulas for approximating $f'(x)$ where $f(x) = \sin(x)$ at $x = \pi/4$. Use $\epsilon = 0.1, 0.01, 0.001$.

\textbf{Hint:} The formulas are:
\begin{itemize}
    \item Forward: $f'(x) \approx \frac{f(x+\epsilon) - f(x)}{\epsilon}$
    \item Backward: $f'(x) \approx \frac{f(x) - f(x-\epsilon)}{\epsilon}$
    \item Central: $f'(x) \approx \frac{f(x+\epsilon) - f(x-\epsilon)}{2\epsilon}$
\end{itemize}
\end{problem}

\begin{problem}[Matrix Conditioning]
\label{prob:matrix-conditioning}
Consider the matrix $\mat{A} = \begin{bmatrix} 1 & 1 \\ 1 & 1.01 \end{bmatrix}$. Compute its condition number and explain why this matrix is ill-conditioned. What happens when you solve $\mat{A}\vect{x} = \vect{b}$ for $\vect{b} = [2, 2.01]$?

\textbf{Hint:} The condition number is $\kappa(\mat{A}) = \|\mat{A}\| \|\mat{A}^{-1}\|$. For 2×2 matrices, you can compute this directly.
\end{problem}

\subsection{Hard Problems (5 problems)}

\begin{problem}[Numerical Stability of Matrix Inversion]
\label{prob:matrix-inversion-stability}
Consider the matrix $\mat{A} = \begin{bmatrix} 1 & 1 \\ 1 & 1 + \epsilon \end{bmatrix}$ where $\epsilon$ is small. Show that the condition number grows as $\epsilon \to 0$. Implement a numerical experiment to demonstrate this and show how the error in $\mat{A}^{-1}$ grows.

\textbf{Hint:} Use the formula $\kappa(\mat{A}) = \frac{\lambda_{\max}}{\lambda_{\min}}$ and compute the eigenvalues analytically.
\end{problem}

\begin{problem}[Gradient Descent Convergence]
\label{prob:gradient-descent-convergence}
Analyze the convergence of gradient descent on the function $f(x, y) = x^2 + 100y^2$ (Rosenbrock function). Derive the optimal learning rate and show that the convergence rate depends on the condition number of the Hessian.

\textbf{Hint:} The Hessian is $\mat{H} = \begin{bmatrix} 2 & 0 \\ 0 & 200 \end{bmatrix}$. The condition number is $\kappa = 100$.
\end{problem}

\begin{problem}[KKT Conditions Application]
\label{prob:kkt-conditions}
Consider the optimization problem:
\begin{align}
\min_{x, y} \quad & x^2 + y^2 \\
\text{subject to} \quad & x + y \geq 1 \\
& x \geq 0, y \geq 0
\end{align}
Find the optimal solution using the KKT conditions and verify that all conditions are satisfied.

\textbf{Hint:} Set up the Lagrangian with multiple constraints and check the complementary slackness conditions.
\end{problem}

\begin{problem}[Mixed Precision Training]
\label{prob:mixed-precision}
Design a numerical experiment to demonstrate the benefits and limitations of mixed precision training. Compare the training dynamics of a simple neural network using FP32 vs FP16 vs mixed precision. Analyze the trade-offs between speed and numerical stability.

\textbf{Hint:} Implement a simple 2-layer network and monitor gradient norms, loss values, and training time.
\end{problem}

\begin{problem}[Numerical Methods for Deep Learning]
\label{prob:numerical-methods-dl}
Implement and compare different numerical methods for computing gradients in a deep neural network:
\begin{itemize}
    \item Automatic differentiation (backpropagation)
    \item Finite differences
    \item Complex-step differentiation
\end{itemize}
Analyze the computational cost and numerical accuracy of each method for networks of different depths.

\textbf{Hint:} Use a simple feedforward network with ReLU activations and compare the gradient computation time and accuracy.
\end{problem}
